El acelerado avance de la inteligencia artificial generativa ha impulsado su adopción en diversos ámbitos de la ingeniería de software, transformando los procesos tradicionales de desarrollo e integración de sistemas interactivos. En la industria de los videojuegos, la generación automatizada de código mediante modelos de lenguaje ofrece optimizar tiempos de producción y facilitar la implementación de mecánicas lógicas complejas. Sin embargo, junto a las capacidades técnicas de estas herramientas, emergen interrogantes fundamentales sobre cómo la procedencia del código influye en la percepción de calidad por parte del usuario final. El presente trabajo de graduación aborda esta problemática delimitando su campo de estudio en la evaluación comparativa entre código desarrollado de forma tradicional y código asistido por inteligencia artificial dentro del motor de videojuegos: Godot; analizando su impacto específicamente en la experiencia del jugador.

Es importante destacar que esta investigación forma parte de un esfuerzo experimental colaborativo que se complementa de manera directa con el trabajo de graduación de la compañera Nicolle Gordillo, titulado «Análisis sobre la percepción de uso de IA en videojuegos mediante pruebas A/B aplicado a trabajo profesional de graduación: Evaluación de la calidad visual y percepción de assets generados con IA en videojuegos pixel art». En este marco de trabajo conjunto, la versión con código e implementación artesanal desarrollada en el presente estudio sirvió como base funcional para la evaluación del apartado gráfico en la investigación complementaria, mientras que el arte en pixel art de corte tradicional presente en los prototipos de esta tesis fue elaborado por Gordillo. Así, mientras que su investigación profundiza en la calidad visual y la percepción estética de los assets gráficos generados con IA, el presente estudio centra su alcance en la experiencia de usuario general y en el rendimiento del código lógico.

La finalidad central de esta investigación es determinar si el uso de herramientas generativas en la programación altera la calidad funcional del software o si las variaciones en la evaluación responden a preconcepciones del usuario sobre la autoría de la tecnología. Para ello, el estudio plantea como hipótesis de investigación ($H_i$) que la diferencia en la experiencia del jugador y en la valoración de los componentes entre un videojuego con código tradicional y uno asistido por inteligencia artificial está significativamente condicionada por el sesgo de etiquetado y las actitudes previas del evaluador hacia la tecnología, frente a una hipótesis nula ($H_0$) que postula la ausencia de diferencias significativas independientemente del conocimiento explícito sobre la autoría.

Para contrastar dicha hipótesis, se diseño una metodología cuantitativa de carácter experimental basada en una prueba a ciegas y una prueba informada (A/B testing) con una muestra multidisciplinaria de 60 estudiantes universitarios. Los participantes evaluaron el rendimiento técnico, la respuesta mecánica y la experiencia de juego utilizando el instrumento estandarizado \textit{Player Experience Inventory} (PXI) junto a métricas que evaluaban el sesgo de etiquetado. Los datos recabados fueron procesados mediante pruebas t de Student para muestras independientes y pareadas, además de análisis de correlación de Pearson, permitiendo aislar la calidad objetiva del código respecto al juicio de percepción subjetivo.

Entre las conclusiones más importantes del trabajo, se determinó que en la condición ciega el código asistido por inteligencia artificial alcanzó una paridad técnica y funcional completa con el desarrollo humano tradicional, no registrandosé variaciones significativas en la experiencia del jugador. No obstante, al revelar la autoría del código para el grupo informado, se observó una caída drástica y estadísticamente significativa en las puntuaciones asignadas (p<0.001), demostrando la presencia de un marcado sesgo de etiquetado. En consecuencia, la evidencia obtenida permitió rechazar la hipótesis nula y aceptar la hipótesis de investigación, concluyendo que la percepción negativa hacia el software asistido por inteligencia artificial no obedece a deficiencias en su arquitecturas o rendimiento, sino a preconcepciones y estigmas actitudinales del usuario. 